\documentclass{article}
\usepackage{bubblecv}

\begin{document}


\begin{cv}[avatar]{Nikunj Arora}{Software Engineer}


\cvsection[summary]{Profile}  %-----------------------------------------------------------

Software engineer experienced in low-level, XR, full-stack, and automation development.


\cvsection[work]{Work experience}  %------------------------------------------------------

\begin{cvevent}[2021][2026]
    \cvname{Specialist Software Engineer}
    \cvdescription{Qt Group, Oulu, Finland}
    \begin{itemize}
        \item Automated \cvhl{Android Automotive} \cvhl{AOSP} image builds via a build server, enabling fast image creation for internal use and client delivery.
        \item Used \cvhl{Android} hidden APIs to enable activity-embedding beyond the public SDK, unlocking custom \cvhl{IVI} UI patterns.
        \item Implemented the \cvhl{Media Sessions API} for \cvhl{Qt for Android Automotive}, enabling media playback control on an \cvhl{IVI}.
        \item Created and maintain the \cvhl{Qt for Android Studio} plugin, letting developers embed \cvhl{QML} in native Android apps without leaving the IDE.
        \item Maintain the companion \cvhl{Qt Gradle plugin} that automates \cvhl{QML} builds inside \cvhl{Android Studio}, removing hand-configured build steps.
        \item Authored \cvhl{Claude}-compatible \cvhl{LLM} skills for \cvhl{Qt for VSCode} to guide new users through onboarding and cut time-to-first-build.
    \end{itemize}
\end{cvevent}

\cvseparator[2]
\begin{cvevent}[2017][2019]
	\cvname{Game Developer}
	\cvdescription{Sameer Hazari Studios Pvt. Ltd., New Delhi, India}
	\begin{itemize}
		\item Built the full-stack \cvhl{AWS} backend and frontend from scratch to underpin the studio's 3D commerce pipeline.
		\item Automated \cvhl{3D} model conversion across formats (incl. \cvhl{USDZ} for Apple devices), replacing manual per-file conversion with a pipeline.
		\item Built an \cvhl{iOS} \cvhl{AR} app and \cvhl{UE4} \cvhl{VR} experience for immersive 3D product previews in e-commerce.
		\item Applied Google's style-transfer \cvhl{deep learning} model to auto-generate product art variations, expanding product catalogs automatically with existing artworks.
		\item Built a browser-side \cvhl{3D} rendering pipeline that produces product videos client-side, removing the need for server-side rendering.
	\end{itemize}
\end{cvevent}


\cvsection[education]{Education}  %------------------------------------------------------

\begin{cvevent}[2019][2021]
    \cvname{Masters in Computer Science}
    \cvdescription{University of Oulu, Finland}
    \begin{itemize}
        \item Specialized in \cvhl{Artificial Intelligence}.
        \item Thesis: Augmenting Virtual Reality telepresence experience using self-avatar, extended to a \cvhl{publication}: \url{https://ieeexplore.ieee.org/document/9714051}
    \end{itemize}
\end{cvevent}


\cvsidebar %-----------------------------------------------------------------------------


\cvsection[contact]{Contact}  %----------------------------------------------------------

\begin{cvitem}[Envelope][4]
    \textbf{Email}\\
    \href{mailto:nikunjarora12345@gmail.com}{\texttt{nikunjarora12345@gmail.com}}
\end{cvitem}

\cvseparator[2]
\begin{cvitem}[Phone][4]
    \textbf{Phone}\\
    \href{tel:+358504795431}{\texttt{+358 50 479 5431}}
\end{cvitem}

\cvseparator[2]
\begin{cvitem}[Globe][4]
    \textbf{Blog}\\
    \href{https://su-nikunj.github.io}{\url{https://su-nikunj.github.io}}
\end{cvitem}

\cvseparator[2]
\begin{cvitem}[Linkedin][4]
    \textbf{Linkedin}\\
    \href{https://linkedin.com/in/nikunj-arora-045067b4}{\url{https://linkedin.com/in/nikunj-arora-045067b4}}
\end{cvitem}

\cvseparator[2]
\begin{cvitem}[CodeFork][4]
    \textbf{Git Repo}\\
    \href{https://github.com/su-nikunj}{\url{https://github.com/su-nikunj}}
\end{cvitem}


\cvsection[skills]{Skills}  %-----------------------------------------------------------

\begin{cvitem}[Code][3]
    \cvhl{Languages}\par
    \cvbadge{C++}\cvbadge{Java}\cvbadge{Python}\par
    \cvbadge{JavaScript}\cvbadge{TypeScript}\cvbadge{Node.js}
\end{cvitem}

\cvseparator[2]
\begin{cvitem}[Cube][3]
    \cvhl{Frameworks}\par
    \cvbadge{Qt}\cvbadge{Android}\cvbadge{iOS}\cvbadge{Unity}\cvbadge{Unreal Engine 4}
\end{cvitem}

\cvseparator[2]
\begin{cvitem}[Wrench][3]
    \cvhl{DevOps \& Tools}\par
    \cvbadge{Linux}\cvbadge{AWS}\cvbadge{Server Management}\cvbadge{Claude code}\cvbadge{Opencode}
\end{cvitem}

\cvsection[person]{References}  %--------------------------------------------------------

\begin{cvitem}
	\cvname{Tinja Paavoseppä}
	\cvdescription{Staff Software Engineer, Qt Group}
	\href{tel: +358505310204}{+358 50 531 0204}\\
	\href{mailto:tinja.paavoseppa@qt.io}{tinja.paavoseppa@qt.io}
\end{cvitem}

\cvseparator[2]
\begin{cvitem}
	\cvname{Nicholas Bennett}
	\cvdescription{Team Lead: AI Tooling, Qt Group}
	\href{tel:+358503231959}{+358 50 323 1959}\\
	\href{mailto:nicholas.bennett@qt.io}{nicholas.bennett@qt.io}
\end{cvitem}

\cvsection[target]{Projects}

\begin{cvitem}
	Self-hosted a local \cvhl{LLM} server with \cvhl{Opencode} as an on-prem dev assistant for privacy and offline use.
\end{cvitem}

\begin{cvitem}
	Migrating my home server to \cvhl{Kubernetes} with a \cvhl{FluxCD} \cvhl{GitOps} workflow, making infrastructure changes version-controlled.
\end{cvitem}


\end{cv}

\end{document}
